% LaTeX file for resume

\documentclass[11pt, a4paper]{article}

\usepackage[top=3.75cm, bottom=3.75cm, left=3.75cm, right=3.75cm]{geometry}
\usepackage{hyperref}
\usepackage{adjustbox}
\usepackage{multicol}
\usepackage{fontspec}
\usepackage{titlesec}
\usepackage{enumitem}
\usepackage{xcolor}
\usepackage{parskip}
\usepackage{microtype}
\usepackage{needspace}

\setlength{\emergencystretch}{5em}
\setlength{\parskip}{10pt plus 2pt minus 1pt}
\hbadness=10000

\hyphenpenalty=10000
\exhyphenpenalty=10000

\setmainfont{Ubuntu}

% Colours
\definecolor{accent}{HTML}{2b2b2b}

% Hyperlink style
\hypersetup{colorlinks=true, urlcolor={blue!60!black}, linkcolor={blue!60!black}}

% Section formatting: small caps, extra space above, no rule
\titleformat{\section}{\large\color{accent}}{}{0em}{}
\titlespacing*{\section}{0pt}{28pt}{10pt}

% Page numbers
\pagestyle{plain}

% Tight lists
\setlist{nosep, leftmargin=0pt, label={}}

\setlength\columnsep{6pt}



\begin{document}

% ---- Header ----
\begin{center}
  {\LARGE\bfseries Thomas Hazell}\\[6pt]
  Somerville College, Woodstock Road, Oxford OX2 6HD\\[2pt]
  \url{thomas.hazell@politics.ox.ac.uk} \enspace$\cdot$\enspace \url{https://t-hazell.github.io/}
\end{center}

\vspace{8pt}

% % ---- Research Interests ----
% \section{Research Interests}
% \vspace*{0pt}
% \begin{multicols}{2}
%   Autocracy\\
%   Bureaucracy \& bureaucrats\\
%   Computational social science\\
%   Central Asia\\
%   Procurement \& corruption\\
%   Subnational governance
% \end{multicols}

\section{Education}

\begin{samepage}
  \textbf{DPhil (PhD) in Politics} \hfill 2024 -- \phantom{2027}\\
  \textit{Department of Politics and International Relations, University of Oxford}

Funded by an Economic and Social Research Council Studentship. Supervised by Dr Katerina Tertytchnaya.
\end{samepage}

\begin{samepage}
  \textbf{Visiting Graduate Student} \hfill March -- April 2026\\
  \textit{Department of Political Science and International Relations, Nazarbayev University, Astana}
\end{samepage}

\begin{samepage}
  \textbf{MPhil in Politics (Comparative Government)} \hfill 2022 -- 2024\\
  \textit{Department of Politics and International Relations, University of Oxford}

Grade: \textit{Distinction}. Funded by an Economic and Social Research Council Studentship. Thesis: `Subnational governance, patronage, and political careers in Kazakhstan,' supervised by Dr Katerina Tertytchnaya.
\end{samepage}

\begin{samepage}
  \textbf{BA in Philosophy, Politics, and Economics} \hfill 2019 -- 2022\\
  \textit{Lincoln College, University of Oxford}

Grade: \textit{First Class}. Finals examinations in politics and economics. Thesis: `Changing the rules of autocracy: the politics of electoral reform in competitive authoritarian Ukraine,' supervised by Dr Jody LaPorte.
\end{samepage}

\section{Work-in-Progress}

For abstracts, see \url{https://t-hazell.github.io/research.html}.

\textbf{Thesis papers}

`Personnel appointments in authoritarian bureaucracies'. \textit{Working paper}.

`Social connections and the quality of public procurement'. \textit{In progress}.

`Which merits matter? Political performance, economic meritocracy, and elite loyalty in authoritarian states'. \textit{Working paper}.

\textbf{Other work}

`The governance effects of local elections in autocracies'. With Kirill Melnikov and Eleonora Minaeva. \textit{Working paper}.

`Local elections and elite management in authoritarian regimes: evidence from Kazakhstan'. With Eleonora Minaeva and Kirill Melnikov.  \textit{Working paper}.

`Shuffling to co-opt: Subnational governance, co-optation, and political careers in Kazakhstan'. \textit{Working paper}.

\section{Invited Talks}

`Personnel appointments in authoritarian bureaucracies.' 20 March 2026. Department of Political Science and International Relations, Nazarbayev University, Astana.

\section{Conferences \& Workshops}

\textbf{Upcoming}

`Personnel appointments in authoritarian bureaucracies.' \textit{American Political Science Association Annual Meeting (APSA)}. 5 Septmber 2026. Boston.

`Elite Selection and Governance in Autocracies: Kazakhstan's Natural Experiment.' (With Kirill Melnikov and Eleonora Minaeva.)  \textit{American Political Science Association Annual Meeting (APSA)}. 4 Septmber 2026. Boston.

`Personnel appointments in authoritarian bureaucracies.' \textit{Conference of the European Political Science Society (EPSS)}. 18 June 2026. Belfast.

`Personnel appointments in authoritarian bureaucracies.' \textit{4th APSG Conference on Authoritarian Politics}. 16 June 2026. Authoritarian Political Systems Group, King's College London.

\textbf{Past}

`Which merits matter? Political performance, economic meritocracy, and elite loyalty in authoritarian states.' \textit{Conference of the European Political Science Association (EPSA)}. 27 June 2025. UC3M, Madrid.

`Personnel appointments in authoritarian bureaucracies.' \textit{3rd APSG Conference on Authoritarian Politics}. 23 June 2025. Authoritarian Political Systems Group, King's College London.

`Personnel appointments in authoritarian bureaucracies.' \textit{3rd Conference ``New Advances in the Political Economy of Development in Eurasia.''} 6 June 2025. Almaty Management University, Almaty.

`Local Elections and Elite Management in Authoritarian Regimes: Evidence from Kazakhstan.' (With Eleonora Minaeva and Kirill Melnikov, presented by Eleonora Mineava.) \textit{Conference of the Midwestern Political Science Association (MPSA)}. 5 April 2025. Chicago.

`Shuffling to co-opt: Subnational governance, patronage, and political careers in Kazakhstan.' \textit{Conference of the Midwestern Political Science Association (MPSA)}. 5 April 2025. Chicago.

`Shuffling to co-opt: Subnational governance, patronage, and political careers in Kazakhstan.' \textit{ESRC Grand Union DTP Annual Conference}. 4 June 2024. University of Oxford, Oxford.

`Shuffling to co-opt: Subnational governance, patronage, and political careers in Kazakhstan.' \textit{Emerging Scholars Workshop on the Study of Challenges in Eurasian Politics and Society}. 24 May 2024. Center for Slavic, Eurasian, and East European Studies, University of North Carolina at Chapel Hill. Presented Online.

\section{Teaching}

\textbf{Undergraduate Tutor} (sole instructor for weekly small-group teaching)\\
\textit{University of Oxford}

Comparative Government (Final exams core paper). TT (Summer) 2026.

Politics in Russia and the Former Soviet Union (Final exams options paper). MT (Autumn) 2025, TT (Summer) 2026.

\section{Professional \& Research Experience}

\begin{samepage}
  \textbf{Research Assistant on Non-Violent Repression in Electoral Autocracies}\hfill{}\linebreak September 2024 -- present\\
  \textit{Department of Politics and International Relations, University of Oxford}

Research assistance to Dr Katerina Tertytchnaya's ESRC-funded project on administrative and bureaucratic repression with an empirical focus on the Russian Federation. Including text-as-data/classification, survey data, and event data tasks.
\end{samepage}

\begin{samepage}
  \textbf{Election Observer (Short-term), Uzbekistan} \hfill October 2024\\
  \textit{OSCE ODIHR Election Observation Mission}

UK-seconded short-term observer for the OSCE ODIHR mission in Samarqand region during the 2024 parliamentary elections.
\end{samepage}

\begin{samepage}
  \textbf{Consultant Researcher} \hfill June -- October 2024\\
  \textit{Ridgeway Information, London}

Consultancy supporting the International Atomic Energy Agency (IAEA)'s nuclear non-proliferation work. Open source methods and Russian-language support, especially on projects related to the former Soviet Union.
\end{samepage}

\begin{samepage}
  \textbf{Research Assistant on AUTHLIB (PI: Stephen Whitefield)} \hfill September 2023 -- January 2024\\
  \textit{Department of Politics and International Relations, University of Oxford}

Part of the joint UKRI/Horizon Europe funded project `neo-authoritarianisms in Europe and the liberal democratic response' (\url{https://www.authlib.eu/}), led at Oxford by Stephen Whitefield, Spyros Kosmidis, and Zofia Stemplowska. Coding tweets from European political parties for NLP training.
\end{samepage}

\begin{samepage}
  \textbf{Research Assistant to Dr Lenka Buštíková} \hfill May -- September 2023\\
  \textit{Oxford School of Global and Area Studies, University of Oxford}

Data collection for projects on illiberalism and irregular armed groups' political role in Ukraine.
\end{samepage}

\begin{samepage}
  \textbf{Research Assistant on the Oxford University Economic Recovery Project} \hfill February -- November 2021\\
  \textit{Smith School of Energy and Enterprise, University of Oxford}

Research assistance, using government and media sources to track state spending in response to Covid-19.
\end{samepage}

\begin{samepage}
  \textbf{Intern} \hfill Summer 2021\\
  \textit{Oxford Internet Institute, University of Oxford}
\end{samepage}

\section{Service}

Reviewer for \textit{Journal of Elections, Public Opinion and Parties} and \textit{Communist and Post-Communist Studies}.

\section{Selected Scholarships \& Awards}

\begin{samepage}
\textbf{Grand Union Doctoral Training Programme Studentship}, Economic and Social Research Council, 2022--2026.

Competitively awarded doctoral studentship, covering fees and stipend for the MPhil in Politics followed by the DPhil (PhD) in Politics, both at Oxford. Value: $\approx$ \pounds 145,000.
\end{samepage}

%\textbf{Language Training Grant}, Department of Politics and International Relations, University of Oxford, 2023--2024.
%
%Funding for a year of Russian courses at the Oxford University Language Centre. Value: \pounds 360.
%
% \begin{samepage}
% \textbf{Research Training Support Grant}, Grand Union Doctoral Training Programme, Summer 2023.

% Grant covering six weeks of intensive Russian language training in Bishkek, Kyrgyzstan. Value: \pounds 820.
% \end{samepage}

\begin{samepage}
\textbf{Tatham Exhibition and College Scholarship}, Lincoln College, University of Oxford, 2020--2022.

Awarded for academic performance. Value: \pounds 450.
\end{samepage}

\section{Languages}

English (Native). Russian (Intermediate). Qazaq (Beginner).

\section{Computer Skills}

\textbf{R} (Advanced): Modelling \& causal inference methods. Survey data. Visualisation. Managing large datasets. Web scraping/automation. Quarto. Machine learning (\texttt{tidymodels}). GIS (\texttt{sf}).

\textbf{Python}: Scraping. Data management. Natural language processing/\allowbreak text-as-data/\allowbreak machine learning, including \texttt{torch}. GIS.

\textbf{\textit{Other}}: OpenAI APIs (Batch, Fine Tuning). Git/GitHub. Cloud computing (Azure, Google Cloud). Databases/\allowbreak basic SQL. \LaTeX.

\section{Additional Training}

\textbf{Qazaq language}, Faculty of Asian and Middle Eastern Studies, University of Oxford, November 2025--February 2026.

\textbf{Comprehensive E-learning Course for OSCE/ODIHR Observers}, September 2024.

\textbf{Fieldwork Safety Overseas}, University of Oxford Safety Office, June 2023.

\textbf{Advanced R for the Social Sciences}, Social Sciences Division, University of Oxford, February 2023.

\section{International Experience}

Intensive Russian language training in Bishkek, Kyrgyzstan, Summer 2022.

Travel in Kyrgyzstan, Kazakhstan, Uzbekistan, Azerbaijan, Georgia, Armenia.

\section{References}

\begin{tabular}[t]{@{}p{0.47\textwidth}@{\hspace{0.06\textwidth}}p{0.47\textwidth}@{}}
  \begin{minipage}[t]{\linewidth}\raggedright
    Dr Jody LaPorte\\
    Gonticas Fellow in Politics and International Relations\\
    Lincoln College\\
    University of Oxford\\
    \url{jody.laporte@lincoln.ox.ac.uk}
  \end{minipage}
  &
  \begin{minipage}[t]{\linewidth}\raggedright
    Dr Katerina Tertytchnaya\\
    Associate Professor of Comparative Politics\\
    Department of Politics and International Relations\\
    University of Oxford\\
    \url{katerina.tertytchnaya@politics.ox.ac.uk}
  \end{minipage}
  % \\[12pt]
  % \begin{minipage}[t]{\linewidth}\raggedright
  %   Professor Lenka Buštíková\\
  %   Department of Political Science\\
  %   University of Florida\\
  %   \url{lenka.bustikova@ufl.edu}
  % \end{minipage}
  % &
\end{tabular}

\end{document}
